\section{Ergebnisse}

\subsection{Merkmale}

In diesem Kapitel werden die Ergebnisse der Feature-Extraktion vorgestellt. Da eine vollständige Beschreibung der Feature-Plots aller elf Merkmale in jeweils 3 Pozessierstufen den Rahmen sprengen würde, wird hier nur auf jeweils eine Prozessierstufe pro Feature eingegangen. Eine vollständige Übersicht aller Merkmale aus allen Prozessierstufen befindet sich im Anhang dieser Arbeit.\\

\subsubsection{Amplitude}

Die Amplituden werden hier in der Prozessierstufe 3 verglichen. Alle Merkmale besitzen hohe Werte in Bereich \textbf{(A)}, welches die Luft- und Bodenwelle zeigt. Desweiteren zeigen sie eine hohe Reflektivität im Bereich \textbf{(B)}, in dem Untergrundstrukturen zum Vorschein kommen und auch Strukturen im Bereich \textbf{(C)} sind überall zu erkennen.

\begin{figure}[H]
    \centering
    \includegraphics[width=\textwidth]{feature_plots_khigh/instantaneous_amplitude.jpg}
    \caption{Instantaneous Amplitude des Radargramms 7, Prozessierstufe 3}
    \label{fig:inst_amp_result}
\end{figure}

\begin{figure}[H]
    \centering
    \includegraphics[width=\textwidth]{feature_plots_khigh/rms_amplitude.jpg}
    \caption{RMS-Amplitude des Radargramms 7, Prozessierstufe 3}
    \label{fig:rms_result}
\end{figure}

\begin{figure}[H]
    \centering
    \includegraphics[width=\textwidth]{feature_plots_khigh/average_energy.jpg}
    \caption{Average Energy des Radargramms 7, Prozessierstufe 3}
    \label{fig:avg_energy_result}
\end{figure}

\begin{figure}[H]
    \centering
    \includegraphics[width=\textwidth]{feature_plots_khigh/quadrature.jpg}
    \caption{Quadrature des Radargramms 7, Prozessierstufe 3}
    \label{fig:quadrature_result}
\end{figure}

Desweiteren wurden alle Merkmale gegeneinander geplottet in einem sogenannten Pairplot. In Abb. \ref{fig:pairplot_amps_result} sieht man die Gegenüberstellungen mit ihren python Variablen Namen. Falls ein Variablenname nicht eindeutig sein sollte, gibt es in der Tabelle \ref{tab:feature_names} eine Übersicht aller Variablen und deren Bedeutung. Die Pairplots zeigen, dass alle Amplitudenmerkmale stark miteinander korrelieren. Besonders auffällig ist die hohe lineare Korrelation zwischen der Average Energy und der Instantaneous Amplitude. RMS Amplitude und Average Energy besitzen fast schon einen quadratischen Zusammenhang, was bei der Berechnung der beiden Feature nicht verwunderlich ist. Rein visuell scheint die Quadrature noch das unabhängigste und gleichverteilteste Merkmal im Vergleich zu den anderen zu sein.

\begin{figure}[H]
    \centering
    \includegraphics[width=\textwidth]{feature_plots_khigh/pairplot_amps.jpg}
    \caption{Pairplot der Amplitudenmerkmale des Radargramms 7, Prozessierstufe 3}
    \label{fig:pairplot_amps_result}
\end{figure}

\subsubsection{Phase}

Die Phase wird in der Prozessierstufe 3 gezeigt. Während die Phasenverlaufe meist sehr gut sichtbar sind, zeigt Bereich \textbf{(A)} Phasenverläufe, die horizontal große Unterschiede aufweisen. Bereich \textbf{(B)} ist dabei eine Außnahme, in dem man die Phasenverläufe bis zur maximalen Laufzeit sehr gut verfolgen kann.

\begin{figure}[H]
    \centering
    \includegraphics[width=\textwidth]{feature_plots_khigh/instantaneous_phase.jpg}
    \caption{InstantaneousPhase des Radargramms 7, Prozessierstufe 3}
    \label{fig:inst_phase_result}
\end{figure}

\subsubsection{Frequenz}

Die Frequenz wird hier zweimal in der Prozessierstufe 1 gezeigt, einmal mit moderater und einmal mit hoher Glättung.

\begin{figure}[H]
    \centering
    \includegraphics[width=\textwidth]{feature_plots_raw/instantaneous_frequency_raw.jpg}
    \caption{Instantaneous Frequency des Radargramms 7, Prozessierstufe 1, moderate Glättung}
    \label{fig:inst_freq_result_1}
\end{figure}

\begin{figure}[H]
    \centering
    \includegraphics[width=\textwidth]{feature_plots_raw/instantaneous_frequency_high.jpg}
    \caption{Instantaneous Frequency des Radargramms 7, Prozessierstufe 1, hohe Glättung}
    \label{fig:inst_freq_result_2}
\end{figure}

\subsubsection{Qualitätsfaktor}

Der Qualitätsfaktor wird in der Prozessierstufe 3 gezeigt. Der Qualitätsfaktor zeigt besonders in Bereich \textbf{(A)} hohe Werte. In Bereich \textbf{(B)} dagegen sind vermehrt Samples mit niedriger Qualität zu erkennen und auch Bereich \textbf{(C)} hebt sich gut vom Hintergrund ab.

\begin{figure}[H]
    \centering
    \includegraphics[width=\textwidth]{feature_plots_khigh/instantaneous_q.png}
    \caption{Qualitätsfaktor des Radargramms 7, Prozessierstufe 3}
    \label{fig:inst_q_result}
\end{figure}

\subsubsection{Statistische Momente}

Die statistischen Momente werden in der Prozessierstufe 3 gezeigt. Die Momente zeigen allgemein zu kleinen Laufzeiten niedrige Werte und zu großen Laufzeiten hohe Werte. In Bereich \textbf{(A)} sind Werte regional etwas kleiner und in Bereich \textbf{(B)} sind die Werte sind die Werte selbst bei langen Laufzeiten noch relativ klein.

\begin{figure}[H]
    \centering
    \includegraphics[width=\textwidth]{feature_plots_khigh/skewness.jpg}
    \caption{Skewness des Radargramms 7, Prozessierstufe 3}
    \label{fig:skewness_result}
\end{figure}

\begin{figure}[H]
    \centering
    \includegraphics[width=\textwidth]{feature_plots_khigh/kurtosis.jpg}
    \caption{Kurtosis des Radargramms 7, Prozessierstufe 3}
    \label{fig:kurtosis_result}
\end{figure}

Nicht nur auf den Feature-Plots selbst sondern auch in ihren Pairplots und in ihrer Werteverteilung sind sich die statistischen Momente sehr ähnlich.

\begin{figure}[H]
    \centering
    \includegraphics[width=\textwidth]{feature_plots_khigh/pairplot_stats.jpg}
    \caption{Pairplot der statistischen Momente des Radargramms 7, Prozessierstufe 3}
    \label{fig:pairplot_statistical_moments_result}
\end{figure}

\subsubsection{Ähnlichkeit}

Die Ähnlichkeit als Semblance wird in Prozessierstufe 2 dargestellt. Die Semblance zeigt in den Bereichen \textbf{(A1 - A3)} eine sehr hohe Ähnlichkeit und in den Bereichen \textbf{(B1 - B3)} eine geringe Ähnlichkeit. Auffällig sind die Bereiche \textbf{(C1 - C3)}, in denen die Semblance horizontale Strukturen von geringer Ähnlichkeit zeigt.

\begin{figure}[H]
    \centering
    \includegraphics[width=\textwidth]{feature_plots/semblance.jpg}
    \caption{Semblance des Radargramms 7, Prozessierstufe 2}
    \label{fig:semblance_result}
\end{figure}

\subsubsection{Neigung}

Der Dip wird ebenfalls in der Prozessierstufe 2 gezeigt. Der Dip zeigt allgemein ein sehr heterogenes Bild. In den Bereichen \textbf{(A1 - A3)} zeigt er eine betragsmäßig hohe Neigung, während die Bereiche \textbf{(B1 - B2)} eine betragsmäßig geringe Neigung aufweisen. In Bereich \textbf{(C)} sind die sonst eher verwaschenen und glatten Übergänge eher scharf.

\begin{figure}[H]
    \centering
    \includegraphics[width=\textwidth]{feature_plots/dip.jpg}
    \caption{Dip des Radargramms 7, Prozessierstufe 2}
    \label{fig:dip_result}
\end{figure}

\subsection{Modellvergleich}

\subsubsection{Startmodell}

Nun werden der Ergebnisse der verschiedenen Testläufe mit verschiedenen Merkmalen und verschiedenen Hyperparametern vorgestellt. Eine trainierte Self-Organizing Map wird als Modell bezeichnet, welches feste Gewichte besitzt. Die ersten Testläufe wurden mit den Merkmalen in Tabelle \ref{tab:model_1_features} durchgeführt. Diese wurden visuell ausgewählt und sollen stark unabhängig voneinander sein.

\begin{table}[H]
    \centering
    \begin{tabular}{|c|c|}
        \hline
        \textbf{Merkmal} & \textbf{Prozessierstufe} \\
        \hline
        Instantaneous Amplitude & 2 \\
        Instantaneous Frequency & 1 \\
        Semblance & 2 \\
        Kurtosis & 2 \\
        Dip & 2 \\
        \hline
    \end{tabular}
    \caption{Merkmale des Modells 1 mit den Hyperparametern $\alpha_0 = 0.3, \sigma_0 = 7, N = 50$}
    \label{tab:model_1_features}
\end{table}

\noindent In Abb. \ref{fig:model_1_planes} ist zu sehen, dass fast alle Merkmale eine gute Sortierung aufweisen. Lediglich die Neigung (Dip) ist nicht gut sortiert, da mehrere kleine Cluster zu sehen sind, aber kein großes zusammenhängendes, wie bei den anderen Merkmalen. Desweiteren fällt auf, dass die vier sortierten Merkmale alle in größtenteils verschiedene Ecken sortiert wurden.

\begin{figure}[H]
    \centering
    \includegraphics[width=\textwidth]{som_runs_pics/R09/som_planes_09_full_on_rg_6.jpg}
    \caption{Model 1: Die jewiligen Werteverteilungen der Neuronen für jedes Merkmal}
    \label{fig:model_1_planes}
\end{figure}

\noindent In der U-Matrix (Abb. \ref{fig:model_1_u_matrix}) sind die durchschnittlichen euklidischen Abstände der Neuronengewichte zu ihren Nachbarn zu sehen. Helle Bereiche zeigen kleine Abstände, also Cluster im Feature-Raum, während dunkle Bereiche große Abstände zeigen, also als Grenzen zwischen Clustern zu verstehen sind. Dabei fällt auf, dass es eine generelle Trennung zwischen weiter oben und weiter unten liegenden Neuronen gibt. Wirft man einen Blick zurück auf Abb. \ref{fig:model_1_planes}, so ist zu erkennen, dass die Neuronen mit hohen Momentanamplituden und hohen Momentanfrequenzen ein größeres Cluster bilden und auch Neuronen mit hohen Semblance- und Kurtosiswerten nah beieinander liegen. Der größte dunkle Bereich, ist auch der Bereich mit hohen Neigungswerten.

\begin{figure}[H]
    \centering
    \includegraphics[width=\textwidth]{som_runs_pics/R09/som_u_matrix_09_full_on_rg_6.jpg}
    \caption{Model 1: Die U-Matrix zeigt die durchschnittliche Distanz der Neuronen im Feature-Raum zu ihren Nachbarn}
    \label{fig:model_1_u_matrix}
\end{figure}

\noindent In der Trefferverteilung (Abb. \ref{fig:model_1_hits}) sind die Trefferanzahl pro Neuron als BMU zu sehen. Dabei ist zu erkennen, dass es größere Cluster in rot, orange, grün, gelb-grün, blau und lila gibt. In Abb. \ref{fig:model_1_bmu} sind die Radargrammpunkte in der Farbe der BMU zu sehen. Dort fallen direkt die blauen Strukturen in der Mitte des Bildes auf, grüne Bereiche im oberen Teil des Bildes und rote Bereiche im unteren Teil des Bildes.

\begin{figure}[H]
    \centering
    \includegraphics[width=\textwidth]{som_runs_pics/R09/som_hits_09_full_on_rg_6.jpg}
    \caption{Model 1: Die Trefferverteilung der Neuronen für jedes Merkmal}
    \label{fig:model_1_hits}
\end{figure}

\begin{figure}[H]
    \centering
    \includegraphics[width=\textwidth]{som_runs_pics/R09/bmu_radargram_09_full_on_rg_6.jpg}
    \caption{Model 1: Die BMU-Verteilung der Neuronen für jedes Merkmal}
    \label{fig:model_1_bmu}
\end{figure}

\subsubsection{Epochenvergleich}

Der Lauf aus dem Startmodell wurden sich auch nach nur 10 und 30 Epochen angeschaut.

\subsubsection{Vergleich verschiedener Lernraten und Nachbarschaftsradien}

Für Modell 2 wurden die gleichen Merkmale verwendet, jedoch mit den Hyperparametern $\alpha_0 = 0.5, \sigma_0 = 5$.

\subsubsection{Vergleich mit anderen Merkmalen}

Nachdem die Neigung wenig sortiert war und dementsprechend wenig zum Endergebnis beigetragen hat, wurde sie in Modell 3 gegen die Merkmale in Tabelle \ref{tab:model_3_features} ersetzt. Die Hyperparameter wurden auf $\alpha_0 = 0.3, \sigma_0 = 7, N = 50$ belassen, um eine gute Vergleichbarkeit mit Modell 1 zu gewährleisten.

\begin{table}[H]
    \centering
    \begin{tabular}{|c|c|}
        \hline
        \textbf{Merkmal} & \textbf{Prozessierstufe} \\
        \hline
        Instantaneous Amplitude & 2 \\
        Instantaneous Frequency & 1 \\
        Semblance & 2 \\
        Kurtosis & 2 \\
        Instantaneous Phase (Real) & 2 \\
        Instantaneous Phase (Imaginary) & 2 \\
        instantaneous Quality Factor & 2 \\
        \hline
    \end{tabular}
    \caption{Merkmale des Modells 3 mit den Hyperparametern $\alpha_0 = 0.3, \sigma_0 = 7, N = 50$}
    \label{tab:model_3_features}
\end{table}

Im Anschluss wurden Features aus der Prozessierstufe 3 verwendet, um Modell 4 zu trainieren. Die benutzen Hyperparameter wurden auf $\alpha_0 = 0.5, \sigma_0 = 5, N = 50$ gesetzt.

\begin{table}[H]
    \centering
    \begin{tabular}{|c|c|}
        \hline
        \textbf{Merkmal} & \textbf{Prozessierstufe} \\
        \hline
        Quadrature & 3 \\
        Instantaneous Frequency & 1 \\
        Kurtosis & 3 \\
        \hline
    \end{tabular}
    \caption{Merkmale des Modells 4 mit den Hyperparametern $\alpha_0 = 0.5, \sigma_0 = 5, N = 50$}
    \label{tab:model_4_features}
\end{table}

\subsubsection{}

\subsection{Weitere Erkenntnisse}

\subsection{Test auf anderem Radargramm}